% !PRINTED 221
\weberpage{221}
% !PRINTED none
% Part-title / section-opening page: large centred "Dritter Theil." heading, subtitle,
% and the "\S. 20" section heading, followed directly by running numbered body text
% (Nr. 196, Nr. 197). No running-head page number is visible in any of the three bands.

\begin{center}
\textbf{Dritter Theil.}\\[4pt]
\textbf{Systeme einzelner und simultaner Formen h\"oheren Grades.}
\end{center}

\begin{center}
\S~20.\\
\gesperrt{Ueberschiebung von Systemen.}
\end{center}

196. \textit{Begriff des Systemes.} Wenn wir eine Anzahl symbolischer Ausdr\"ucke in ihrer Gesammtheit bezeichnen wollen, so werden wir dasselbe kurz mit dem Namen \glqq System von Formen\grqq{} belegen. Der Begriff eines Formensystems ist damit in seinen weitesten Grenzen gegeben. Die Formen desselben k\"onnen einzelne symbolische Producte sein mit oder ohne Zahlenfactoren. Sie k\"onnen aber auch ganze Aggregate symbolischer Producte repr\"asentiren; nur wollen wir in diesem Falle voraussetzen, dass alle Glieder des Aggregates die gleiche Anzahl von Symbolen der etwa auftretenden Formen enthalten. Das ganze System kann alle Covarianten und Invarianten eines simultanen Systems bin\"arer Formen $f_1, f_2, f_3 \ldots f_n$ umfassen; es kann sich aber auch auf irgend eine durch gewisse Eigenschaften definirte Gruppe derselben beschr\"anken. Dabei ist nicht nothwendig, dass alle Co- und Invarianten eines solchen Systems von einander unabh\"angig sind; doch werden wir im allgemeinen wenigstens die Annahme machen, dass keine eine \gesperrt{ganze} Function von andern Formen des Systemes ist.

197. \textit{Ueberschiebung zweier Systeme.} Denken wir uns nun zwei solche Formensysteme mit je einer endlichen Anzahl von Formen gegeben, einmal das System
\weq{I}{\overset{a_1}{A_1(x)}, \quad \overset{a_2}{A_2(x)}, \quad \overset{a_3}{A_3(x)} \ldots \overset{a_r}{A_r(x)},}
sodann ein zweites System
\weq{II}{\overset{b_1}{B_1(x)}, \quad \overset{b_2}{B_2(x)}, \quad \overset{b_3}{B_3(x)} \ldots \overset{b_s}{B_s(x)},}
wobei mit $A_i(x)$, $B_i(x)$ irgend welche symbolische Producte oder ganze Aggregate von solchen bezeichnet sind, und die dar\"uber geschriebenen Zahlen $a_i$, $b_i$ ihren Grad in $x$ ausdr\"ucken.


% ===== pdf idx 236, printed 222.0 =====
% !PRINTED 222
\weberpage{222}
Wir stellen uns die Aufgabe, beide Systeme \"ubereinander zu schieben. Wie wir fr\"uher die Ueberschiebung zweier Formen dadurch bildeten, dass wir an dem Product derselben alle m\"oglichen Faltungen des betreffenden Ueberschiebungsgrades ausf\"uhrten, so werden wir auch jetzt eine Ueberschiebung der beiden Systeme erhalten, indem wir irgend ein Product aus den Formen des Systemes (I) mit einem Product von Formen des Systemes (II) falten.

Sei
\[
U = A_1{}^{\alpha_1} A_2{}^{\alpha_2} \ldots A_r{}^{\alpha_r}
\]
ein Product aus den Formen des ersten,
\[
V = B_1{}^{\beta_1} B_2{}^{\beta_2} \ldots B_s{}^{\beta_s}
\]
ein Product aus den Formen des zweiten Systemes, wobei die Exponenten $\alpha_i$, $\beta_i$ irgend welche ganze Zahlen, Null nicht ausgeschlossen, bedeuten.

Das Product $U$ ist in $x$ vom Grade
\[
a_1 \alpha_1 + a_2 \alpha_2 + \cdots a_r \alpha_r;
\]
das Product $V$ vom Grade
\[
b_1 \beta_1 + b_2 \beta_2 + \cdots b_s \beta_s.
\]
Um eine beliebige Ueberschiebung $(U, V)^{\nu}$ zu bilden, haben wir (vgl. Nr. 38) das Product $U \cdot V$ auf alle m\"oglichen Arten $\nu$ mal zu falten und die erhaltene Summe von symbolischen Producten durch die Anzahl ihrer Glieder zu dividiren. Jedes dieser Glieder besitzt einmal die \gesperrt{alten}, von vornherein dem ganzen Aggregat gemeinschaftlichen Klammerfactoren, die ja an der Faltung nicht betheiligt sind, dann aber auch \gesperrt{$\nu$ neue} Klammerfactoren, in deren jedem ein Symbol des Systemes der $A_i(x)$ mit einem Symbol des Systemes der $B_i(x)$ verbunden ist; ausserdem enthalten diese Glieder im Allgemeinen noch Factoren \gesperrt{erster Art}. Ist der Grad eines Gliedes der Ueberschiebung in den Factoren erster Art aus dem Systeme (I) noch $\varrho$, und ebenso in den Factoren erster Art aus dem Systeme (II) noch $\sigma$, dann m\"ussen f\"ur jedes Glied und damit auch f\"ur die ganze Ueberschiebung $(U, V)^{\nu}$ die Gleichungen bestehen:
\weq{1}{a_1 \alpha_1 + a_2 \alpha_2 + a_3 \alpha_3 \cdots + a_r \alpha_r = \varrho + \nu}
\weq{2}{b_1 \beta_1 + b_2 \beta_2 + b_3 \beta_3 \cdots + b_s \beta_s = \sigma + \nu.}
Ist umgekehrt ein Zahlensystem $\alpha_i$, $\beta_i$, $\nu$ gegeben, das diese Gleichungen befriedigt, so k\"onnen wir uns auch dazu eine Ueberschiebung der Producte $U$ und $V$ verschaffen. Denn durch die Exponenten $\alpha_i$ ist das Product $U$, durch die Exponenten $\beta_i$ das Product
%CONT>>


% ===== pdf idx 237, printed 223.0 =====
% !PRINTED 223
\weberpage{223}
%<<CONT
$V$ eindeutig bestimmt, und zwei solche Producte haben nur Eine Ueberschiebung vom Grade $\nu$.

198. \textit{Umgrenzung der zu betrachtenden Ueberschiebungen.} Die beiden Systeme k\"onnen nat\"urlich auf diese Weise zu unz\"ahligen Ueberschiebungen Veranlassung geben. Denn die beiden Gleichungen (1) und (2) sind lineare diophantische Gleichungen mit den Unbekannten $\alpha_i$, $\beta_i$, $\nu$, und zwar solche (vgl. Bd.\ I, Nr. 187), die immer positive ganzzahlige L\"osungen besitzen. Jeder Ueberschiebung entspricht aber eine solche L\"osung, und umgekehrt sahen wir soeben, geh\"ort jeder positiven ganzzahligen L\"osung eine und nur eine Ueberschiebung zu. Von dieser unendlich grossen Anzahl von Ueberschiebungen interessiren uns aber nun jene, durch welche wesentlich neue Formen entstehen. Es werden n\"amlich wie fr\"uher bei der Systembildung einfacher Formen die Ueberschiebungen in zwei Gruppen zu theilen sein: in eine erste Gruppe von Ueberschiebungen, bei denen die Faltung des Productes $U \cdot V$ kein zerfallendes Glied liefert, und in eine zweite, welche solche Glieder enthalten. Hierbei nennen wir ein Glied $G$ ein \gesperrt{zerfallendes}, wenn man dasselbe in zwei Theile $G_1$ und $G_2$ spalten kann, von denen der eine, also auch der andere, ein Glied einer fr\"uheren, niedrigeren Ueberschiebung ist, oder auch ein Product urspr\"unglicher Formen. Die Ueberschiebungen der zweiten Gruppe aber k\"onnen keine wesentlich \gesperrt{neuen} Formen liefern, da sich ja nach den allgemeinen S\"atzen in \S~3 jede solche Ueberschiebung durch dieses Glied und niedrigere Ueberschiebungen, also durch Bekanntes ausdr\"ucken l\"asst. (Siehe als Beisp. etwa Nr. 211.) Wir beschr\"anken uns also auf das Studium der ersten Ueberschiebungsgruppe, deren Formen wir in diesem Paragraphen mit $C_i$ bezeichnen werden. Von diesen Formen werden wir in der Folge vier Lehrs\"atze aufstellen, mit deren Hilfe die Endlichkeit des Formensystems einer bin\"aren Form sich einfach beweisen l\"asst. Hiebei bezeichnen wir als Form $C_i$ nur ein einziges \"ubrigens willk\"urliches Glied der betreffenden Ueberschiebung, die nach unsern allgemeinen Ueberschiebungss\"atzen durch ein solches Glied vollst\"andig definirt ist.

199. \textit{Erster Lehrsatz:} \gesperrt{\glqq Ist das System der Formen $A_i$ sowohl als auch das System der Formen $B_i$ ein endliches, so ist auch das durch Ueberschiebung entstandene System der Formen $C_i$ endlich.\grqq{}}

Der Beweis wird mit Hilfe des im ersten Bande dieser Vorlesungen \S~15 bewiesenen Satzes erbracht, wonach ein simultanes System linearer diophantischer Gleichungen nur eine endliche Zahl einfacher L\"osungen besitzt. Wir sahen n\"amlich soeben, dass jeder
%CONT>>


% ===== pdf idx 238, printed 224.0 =====
% !PRINTED 224
\weberpage{224}
%<<CONT
L\"osung der beiden diophantischen Gleichungen (1) und (2) Nr. 197 eine und nur eine Ueberschiebung entspricht. K\"onnen wir also zeigen, dass jeder zusammengesetzten L\"osung eine Ueberschiebung mit zerfallenden Gliedern zugeh\"ort und umgekehrt, dann k\"onnen die Ueberschiebungen mit nicht zerfallenden Gliedern, also unsere Formen $C_i$, nur den einfachen L\"osungen entsprechen. Diese sind aber nur in endlicher Zahl vorhanden, also auch die Formen $C_i$.

Angenommen aber, die Ueberschiebung $(U, V)^{\nu}$ enthalte ein Glied
\[
G = G_1 \cdot G_2,
\]
wobei $G_1$ und $G_2$ Glieder niedrigerer Ueberschiebungen vom Grade $\nu_1$ bezw. $\nu_2$ sind, so dass $\nu = \nu_1 + \nu_2$, dann l\"asst sich zu jedem der beiden Glieder je eine L\"osung $L_1 = (\alpha_i^{(1)}, \beta_i^{(1)}, \nu_1)$, beziehungsweise $L_2 = (\alpha_i^{(2)}, \beta_i^{(2)}, \nu_2)$ angeben, welche diese niedrigeren Ueberschiebungen v\"ollig bestimmen. Dabei ist alsdann, wenn $L = (\alpha_i, \beta_i, \nu)$ die L\"osung ist, durch welche die Ueberschiebung mit dem Gliede $G = G_1 \cdot G_2$ erhalten wurde:
\[
\begin{aligned}
\nu' &= \nu_1 + \nu_2 \\
\alpha_i &= \alpha_i^{(1)} + \alpha_i^{(2)} \\
\beta_i &= \beta_i^{(1)} + \beta_i^{(2)},
\end{aligned}
\]
d.\,h. die L\"osung $L$ l\"asst sich zusammensetzen aus den L\"osungen $L_1$ und $L_2$. Umgekehrt: Entsprechen den L\"osungen $L_1$ und $L_2$ bezw. die Ueberschiebungen $(U_1, V_1)^{\nu_1}$ und $(U_2, V_2)^{\nu_2}$, so entspricht der zusammengesetzten L\"osung $L = L_1 + L_2$ nothwendig die Ueberschiebung:
\[
(U_1 U_2, \ V_1 V_2)^{\nu_1+\nu_2}
\]
und diese enth\"alt zerfallende Glieder. Denn unter den Gliedern, die aus dem Producte
\[
U_1\, U_2\, V_1\, V_2
\]
hervorgehen, indem man dasselbe auf alle m\"oglichen Arten $(\nu_1+\nu_2)$-mal faltet, treten auch Producte von Gliedern auf, die entstehen, wenn man erst $U_1 V_1$ $\nu_1$-mal und dann $U_2 V_2$ $\nu_2$-mal faltet, die also Glieder der niedrigeren Ueberschiebungen
\[
(U_1, V_1)^{\nu_1} \text{ und } (U_2, V_2)^{\nu_2}
\]
darstellen. Es k\"onnen also in der That die Ueberschiebungen \gesperrt{ohne zerfallende Glieder nur den einfachen L\"osungen entsprechen}, und deren Anzahl ist eine endliche.

200. \textit{Definition eines vollst\"andigen Systemes.} Wir nennen nun ein System von Formen $A_i$ ein \gesperrt{vollst\"andiges}, sobald aus einer einzelnen Form oder einem Producte derselben durch Faltung keine wesentlich neue Form entstehen kann, sondern immer nur Formen, welche ganze rationale Functionen der $A_i$ sind.


% ===== pdf idx 239, printed 225.0 =====
% !PRINTED 225
\weberpage{225}
Solche vollständige Systeme sind die bereits früher aufgestellten
Systeme der Form zweiter, dritter oder vierter Ordnung. Jede
Ueberschiebung $(U,\ U')^\nu$ zweier Producte $U$ und $U'$ der Formen
$A_i$ ist also darstellbar in der Form
\[ (U,\ U')^\nu = \boldsymbol{F}(A_1,\ A_2,\ \ldots\ A_r) = \boldsymbol{F}(A_i), \]
wo $\boldsymbol{F}$ eine ganze rationale Function der Formen $A_i$ bedeutet. So ist
z. B. für die Form vierter Ordnung
\[ (f^{\alpha_1}\cdot \Delta^{\alpha_2}\cdot t^{\alpha_3}\cdot i^{\alpha_4}\cdot j^{\alpha_5},\ \ f^{\beta_1}\cdot \Delta^{\beta_2}\cdot t^{\beta_3}\cdot i^{\beta_4}\cdot j^{\beta_5})^\nu = \boldsymbol{F}(f,\ \Delta,\ t,\ i,\ j). \]

201. \textit{Zweiter Lehrsatz.} Sind die Systeme $(A)$ und $(B)$
vollständige Formensysteme mit den Formen
\weq{I}{ A_1,\ A_2,\ A_3,\ \ldots\ A_r }
\weq{II}{ B_1,\ B_2,\ B_3,\ \ldots\ B_s, }
gelten also für diese die Bedingungen
\weq{1}{ (U,\ U')^\nu = \boldsymbol{F}(A_i) }
\weq{2}{ (V,\ V')^\nu = \boldsymbol{\Phi}(B_i), }
dann führt wiederum die Ueberschiebung derselben auf ein drittes
System selbstständiger Formen
\weq{III}{ C_1,\ C_2,\ C_3,\ \ldots\ C_p, }
zu welchem selbstverständlich als nullte Ueberschiebungen auch die
Formen der beiden ersten Systeme gehören.

Von diesem System können wir behaupten:

\glqq Das System der Formen $C_i$ ist ein vollständiges, sobald die
ursprünglichen Systeme $(A)$ und $(B)$, aus denen dasselbe durch
Ueberschiebung entstanden ist, vollständig sind.\grqq{}

Unsere Aufgabe besteht darin, den Nachweis zu führen, dass auch
für das System der Formen $C_i$, wie für die beiden ursprünglichen
Systeme, die Gleichung besteht
\weq{3}{ (\Gamma,\ \Gamma')^\nu = \boldsymbol{\Psi}(C_i), }
wenn $\Gamma$ und $\Gamma'$ beliebige Producte der $C_i$ sind.

Nun entsteht jede dieser Ueberschiebungen (3) durch Faltung der
in $\Gamma$ enthaltenen Formen $C_i$ mit jenen, die sich in $\Gamma'$
befinden. Diese Formen $C_i$ sind aber aus Producten $U_i$ der Formen
$A_i$ durch Faltung mit Producten $V_i$ der Formen $B_i$ hervorgegangen.
Also kann man sich auch die Form $\boldsymbol{\Psi}$ direct durch Faltung eines
Productes $U\cdot V$ entstanden denken, wobei einmal Faltungen an den
Factoren $A_i$ von $U$, sodann Faltungen an den Factoren $B_i$ von $V$
und endlich Faltungen der Factoren $A_i$ mit Factoren $B_i$ auszuführen
sind. Bezeichnen wir das durch Faltungen der ersten Art aus $U$
entstehende Product mit $\overline U$, das durch Faltungen der zweiten
Art aus $V$ entstehende mit
%CONT>>


% ===== pdf idx 240, printed 226.0 =====
% !PRINTED 226
\weberpage{226}
%<<CONT
$\overline V$, dann entsteht $\boldsymbol{\Psi}$ aus dem Producte $\overline U\cdot\overline V$
durch Faltungen der dritten Art, und kann somit als Glied der
Ueberschiebung
\[ (\overline U,\ \overline V)^\varrho,\quad \varrho \leqq \nu, \]
angesehen werden. Daher lässt sich $\boldsymbol{\Psi}$ durch diese und niedrigere
Ueberschiebungen dieser Art nach \S~3 Nr. 44 darstellen, d. h. es ist
\[ (\Gamma,\ \Gamma')^\nu = \sum (\overline U,\ \overline V)^{\varrho_i}. \]
Da aber $\overline U$ durch Faltung von Formen $A_i$, $\overline V$ durch Faltung
von Formen $B_i$ entstanden ist, so kann nach Voraussetzung
\weq{4}{ \left.\begin{aligned} \overline U &= \sum (U'',\ U')^\nu = \boldsymbol{F}(A_i) \\ \overline V &= \sum (V'',\ V')^\mu = \boldsymbol{\Phi}(B_i) \end{aligned}\right\} }
gesetzt werden, und demnach ist
\[ (\Gamma,\ \Gamma')^\nu = \sum \big(\boldsymbol{F}(A_i),\ \boldsymbol{\Phi}(B_i)\big)^{\varrho_i}; \]
weil aber die Functionen $\boldsymbol{F}$ und $\boldsymbol{\Phi}$ rational und ganz in den
Formen $A_i$ bezw. $B_i$ sind, also
\[ \boldsymbol{F} = \sum A_1^{\alpha_1} A_2^{\alpha_2} \ldots A_r^{\alpha_r} = \sum U \]
\[ \boldsymbol{\Phi} = \sum B_1^{\beta_1} B_2^{\beta_2} \ldots B_s^{\beta_s} = \sum V, \]
und weil die Ueberschiebung von Summen gleich der Summe der
Ueberschiebungen ihrer Posten ist, so wird endlich
\weq{5}{ (\Gamma,\ \Gamma')^\nu = \sum (U,\ V)^{\varrho_i}. }
Die Ueberschiebungen $(U,\ V)^{\varrho_i}$ sind aber nun entweder
Ueberschiebungen der ersten Gruppe, d. h. (vergl. Nr. 198) sie enthalten
keine zerfallenden Glieder und wir können sie direct durch die Formen
$C_i$ ersetzen, oder aber sie enthalten zerfallende Glieder
$G_i = G_1\cdot G_2$; dann ersetzen wir sie zunächst nach Nr. 44 durch
$G_1\cdot G_2$ plus Glieder \gesperrt{niedrigerer Ueberschiebungen}
$(U,\ V)^\lambda$, wo $\lambda < \varrho_i$. Gleichzeitig führen wir alsdann
auch für die Factoren $G_1$ wie $G_2$ die ihnen einzeln entsprechenden
niedrigeren Ueberschiebungen nach demselben Satze Nr. 44 ein. Damit
ist eine erste Reduction der rechten Seite der Gleichung hergestellt.
Diese Reduction werden wir nöthigenfalls so lange fortsetzen, bis nur
mehr nullte Ueberschiebungen, d. i. Producte der $C_i$ auftreten. Die
rechte Seite der letzten Gleichung lässt sich also in der That auf
eine Function der Formen $C_i$ successive reduciren und die
Behauptung:
\[ (\Gamma,\ \Gamma')^\nu = \boldsymbol{\Psi}(C_i) \]
ist sonach erwiesen.


% ===== pdf idx 241, printed 227.0 =====
% !PRINTED 227
\weberpage{227}
202. \textit{Definition relativ vollständiger Systeme.} Es kann nun
aber auch der Fall eintreten, dass durch Faltung des Productes $U$
der Formen $A_i$, die einem bestimmten Systeme (I) angehören, Formen
$\overline U$ entstehen, welche der Gleichung
\weq{1}{ \overline U = \boldsymbol{F}(A_i) + \sum G\cdot W }
genügen. Dabei ist $\boldsymbol{F}(A_i)$ wieder eine ganze rationale Function
der $A_i$, $\sum G\cdot W$ aber ein Aggregat von symbolischen Producten,
deren jedes mindestens einen aus mehreren von vornherein bezeichneten
Klammerfactoren $G$ (oder auch ein Product solcher) zum Factor
besitzt. Diese Gruppe von Klammerfactoren $G$ nennen wir \glqq
Moduln\grqq{} (vergl. Nr. 43) des Formensystemes $(A)$, und sagen
alsdann:

\glqq Das Formensystem $(A)$ ist relativ vollständig in Bezug auf
die Moduln $G$.\grqq{}

Wir drücken dies durch die Gleichung (1) aus oder auch durch die
Congruenz:
\weq{2}{ \overline U \equiv \boldsymbol{F}(A_i)\ \mathrm{mod.}\ G. }
So bilden z. B. im Formensystem der Form vierter Ordnung $f=a_x^4$
die Formen
\[ f = a_x^4,\qquad t = (ab)^2(ca)a_xb_x^2c_x^3 \]
\[ \Delta = (ab)^2 a_x^2 b_x^2,\quad j = (ab)^2(bc)^2(ca)^2 \]
für sich ein relativ vollständiges System modulo $(ab)^4$, insofern
alle aus dem Product
\[ f^\lambda\cdot \Delta^\mu\cdot t^\nu\cdot j^\sigma = \Gamma \]
durch Faltung entstehenden Formen ganze Functionen dieser vier
Formen sind bis auf Glieder mit dem Factor $(ab)^4$. Wir drücken
dies aus durch
%CONT>>


% ===== pdf idx 242, printed 228.0 =====
% !PRINTED 228
\weberpage{228}
%<<CONT
\[ \overline\Gamma = (\Gamma,\ \Gamma')^\nu = \boldsymbol{F}(f,\ \Delta,\ t,\ j) + \sum (ab)^4 W, \]
oder
\[ \overline\Gamma = (\Gamma,\ \Gamma')^\nu \equiv \boldsymbol{F}(f,\ \Delta,\ t,\ j)\ \mathrm{mod.}\ (ab)^4.\footnote{Es dürfte hiebei die selbstverständliche Bemerkung erlaubt sein, dass, wenn $i=(ab)^4$ identisch null ist, dieses relativ vollständige System zu einem absolut vollständigen wird. (Vergl. auch Nr. 209.)} \]

203. \textit{Dritter Lehrsatz.} Sind nun zwei Systeme $(A)$ und $(B)$
gegeben, deren jedes relativ vollständig mod. $G$ ist, dann können
wir wiederum durch Ueberschiebung derselben ein drittes System $(C)$
erhalten. Von diesem System gilt der Satz:

\glqq Das System der Formen $C_i$ ist relativ vollständig in Bezug
auf die Moduln $G$, sobald es auch die beiden ursprünglichen
Systeme $(A)$ und $(B)$ sind, aus denen das System $(C)$ durch
Ueberschiebung entstanden ist.\grqq{}

Wir haben wiederum, wenn $\Gamma$ und $\Gamma'$ zwei beliebige
Producte der Formen $C_i$ sind, die Richtigkeit der Gleichung:
\[ (\Gamma,\ \Gamma')^\nu = \boldsymbol{\Psi}(C_i) + \sum W\cdot G, \]
nachzuweisen. Der Beweis gestaltet sich analog wie der des zweiten
Lehrsatzes. Zunächst schliessen wir ganz in derselben Weise wie
dort, dass sich die Form $(\Gamma,\ \Gamma')^\nu$ als Glied der
Ueberschiebung
\[ (\overline U,\ \overline V)^\varrho,\quad \varrho \leqq \nu, \]
darstellen lässt, wobei $\overline U$ und $\overline V$ Formen sind, die
aus $U$ resp. $V$ durch Faltung der Formen $A_i$ resp. $B_i$
entstehen. Es lässt sich demnach $(\Gamma,\ \Gamma')^\nu$ wiederum
durch eine Summe von Ueberschiebungen vom Typus $(\overline U,\ \overline V)^\varrho$
darstellen, so dass man hat
\[ (\Gamma,\ \Gamma')^\nu = \sum (\overline U,\ \overline V)^{\varrho_i}, \]
Nun ist aber nach Voraussetzung
\[ \overline U = \boldsymbol{F}(A_i) + \sum G\cdot W \]
\[ \overline V = \boldsymbol{\Phi}(B_i) + \sum G\cdot W', \]
also wird
\[ (\Gamma,\ \Gamma')^\nu = \sum \big(\boldsymbol{F}(A_i) + \sum G\cdot W,\ \ \boldsymbol{\Phi}(B_i)+\sum G\cdot W'\big)^{\varrho_i}, \]
oder nach Nr. 32, $(\beta)$
\[ = \sum \big(\boldsymbol{F}(A_i),\ \boldsymbol{\Phi}(B_i)\big)^{\varrho_i} \ + \sum \big(\boldsymbol{F}(A_i),\ \sum G\cdot W'\big)^{\varrho_i} \]
\[ {}+ \sum \big(\sum G\cdot W,\ \boldsymbol{\Phi}(B_i)\big)^{\varrho_i} + \sum \big(\sum G\cdot W,\ \sum G\cdot W'\big)^{\varrho_i}. \]

Der erste Term rechts reducirt sich in derselben Weise wie früher
auf $\sum (U,\ V)^{\varrho_i}$; die übrigen Terme auf Formen, welche die
Moduln $G$, die ja als Klammerfactoren nicht gefaltet werden können,
zu Factoren haben. Daher können wir die letzte Gleichung schreiben
\[ (\Gamma,\ \Gamma')^\nu = \sum (U,\ V)^{\varrho_i} + \sum G\cdot Q. \]

Die Ueberschiebungen $(U,\ V)^{\varrho_i}$ führen aber entweder wieder
auf Formen $C_i$ oder sind Ueberschiebungen mit zerfallenden
Gliedern $G_i$ und wir können alsdann genau wie in Nr. 201 die
allmähliche Reduction auf Formen $C_i$ vornehmen. Also ist in der
That
\[ (\Gamma,\ \Gamma')^\nu \equiv \boldsymbol{\Psi}(C_i)\ \mathrm{mod.}\ G. \]


% ===== pdf idx 243, printed 229.0 =====
% !PRINTED 229
\weberpage{229}
204. \textit{Vierter Lehrsatz.} Von den Moduln, in Bezug auf welche ein
erstes System (A) relativ vollständig ist, können einige als Klammerfactoren
von Formen $B_i$ eines zweiten Systemes (B) auftreten. Wir wollen diese
Moduln mit $H$, die übrigen des Systemes (A) mit $G$ bezeichnen. Es möge
nun das System (A) der Gleichung genügen:
\weq{1}{\bar U = F(A_i) + \sum H\cdot W + \sum G\cdot W'.}

Nehmen wir alsdann an, dass im System (B), das wir gleichfalls relativ
vollständig mod. $G$ voraussetzen, zu jedem der Moduln $H$ eine
entsprechende Form $B_i$ vorhanden sei, die neben Factoren erster Art nur
diesen Modul als Klammerfactor hat, dann genügt System (B) immer noch der
Gleichung
\weq{2}{\bar V = \Phi(B_i) + \sum G\cdot W'',}
wobei allerdings unter den Klammerfactoren von $W''$ auch Moduln $H$
vorhanden sein können, die wir jedoch hier neben den Moduln $G$ nicht
weiter hervorheben wollen.

Durch Ueberschiebung beider Systeme (A) und (B) entsteht ein drittes
System (C). Von ihm gilt der Satz:

\glqq Genügt das System (A) der Gleichung (1), das System (B) aber, das
je eine Covariante besitzen möge, die je einem Modul $H$ von (A)
entspricht, der Gleichung (2), so genügt auch das System (C) der
Gleichung
\[ (\Gamma,\ \Gamma')^\nu = F(C_i) + \sum G\cdot W''', \]
d. h. es ist relativ vollständig modulo $G$.\grqq{}

Ehe ich den Beweis antrete, will ich den Fall durch ein Beispiel
illustriren. Im System der Form $f=a_x^6$ bilden, wie wir demnächst sehen
werden, die drei Formen
\[ f = a_x^6,\quad H = (ab)^2 a_x^2 b_x^2,\quad t = (ab)^2(ca) a_x^3 b_x^4 c_x^5 \]
für sich ein in Bezug auf die beiden Moduln $(ab)^4$ und $(ab)^6$ relativ
vollständiges System (A), d. h. es ist:
\[ \bar U = F(f,\ H,\ t) + \sum (ab)^4\cdot W + \sum (ab)^6\cdot W'. \]

Stellen wir diesem System (A) das volle System (B) der biquadratischen
Covariante $k=(f,f)^4=(ab)^4a_x^2b_x^2$ von $f$ gegenüber
\weq{B}{\left.\begin{aligned}
B_1&=k, & B_2&=(k,\ k)^2, & B_3&=(k,\ k)^4\\
B_4&=(k,\ (k,\ k)^2), & B_5&=((k,\ k)^2,\ k)^2
\end{aligned}\right\},}
so ist dasselbe gerade ein System mit einer Form (nämlich der Grundform
$B_1=k$ selbst), die den Modul $(ab)^4$ besitzt, und das im übrigen,
%CONT>>


% ===== pdf idx 244, printed 230.0 =====
% !PRINTED 230
\weberpage{230}
%<<CONT
wie der nächste Paragraph schon lehren wird, relativ vollständig ist
modulo $(ab)^6$. Durch Ueberschiebung beider Systeme entsteht das System
(C), welches der Gleichung genügt:
\[ (\Gamma,\ \Gamma')^\nu = F(C_i) + \sum (ab)^6\cdot W'''. \]

Wir werden sehen, dass dieses System (C) zusammen mit der Invariante
$(ab)^6$ selbst, gerade das volle System von $f=a_x^6$ bildet.

Der Beweis unseres Satzes gestaltet sich nun wie die beiden
vorhergehenden. Es ist wiederum zunächst:
\[
\begin{aligned}
(\Gamma,\ \Gamma')^\nu &= \sum (\bar U,\ \bar V)^{\varrho_i}, \quad \varrho_i \leqq \nu\\
&= \sum \bigl(F(A_i) + \sum H\cdot W + \sum G\cdot W',\ \ \Phi(B_i) + \sum G\cdot W''\bigr)^{\varrho_i},
\end{aligned}
\]
oder:
\weq{3}{(\Gamma,\ \Gamma')^\nu = \sum \bigl(F(A_i),\ \Phi(B_i)\bigr)^{\varrho_i} + \sum \bigl(H\cdot W,\ \Phi(B_i)\bigr)^{\varrho_i} + \sum G\cdot Q.}

Hierin führt der erste Term rechts in bekannter Weise auf Formen $C_i$;
der dritte besitzt nur Formen, welche Moduln $G$ zu Factoren haben. Es
erübrigt also nur die Discussion des zweiten Termes. Nun kann $H\cdot W$
als Glied der Ueberschiebung $(B_i,\ U_1)^\lambda$ erhalten werden, eben
weil der Modul $H$ auch einer Form $B_i$ als Klammerfactor angehört
(vergl. Nr. 45), und lässt sich folglich durch solche Ueberschiebungen
darstellen. Man kann daher statt des zweiten Termes auch schreiben:
\[ \sum \sum \bigl((B_i,\ \bar U_1)^{\lambda_k},\ \Phi(B_i)\bigr)^{\varrho_i}. \]

Die Glieder dieser Ueberschiebungen entstehen durch Faltung aus Producten
\[ B_i \cdot \bar U_1 \cdot \Phi(B_i), \]
also aus Producten, die eine Form $B_i$ mehr als das ursprüngliche $\bar
V$ und demnach eine Form $A_i$ weniger als das ursprüngliche Product
$\bar U$ in der Gleichung $(\Gamma,\ \Gamma')^\nu = \sum (\bar U,\ \bar
V)^{\varrho_i}$ enthalten. Ersetzen wir nun das gegenüber $\bar U$
reducirte Product $\bar U_1$ durch seinen Werth $F_1(A_i) + \sum H\cdot
W_1 + \sum G\cdot W_1'$, wobei in jedem dieser drei Terme eine Form $A_i$
weniger betheiligt ist als in (1), so gelangen wir genau durch dieselben
Schlüsse wie bisher zu einer abermaligen Reduction des zweiten Termes in
(3), indem derselbe hiedurch auf Producte zurückgeführt wird, die aus
\[ B_k \cdot B_i \cdot \bar U_2 \cdot \Phi(B_i) \]
durch Faltung entstehen. Hierin besitzt das Product $\bar U_2$ abermals
einen Factor $A_i$ weniger als $\bar U$. Durch fortgesetzte Wiederholung
dieses Verfahrens wird schliesslich der Ausdruck
%CONT>>


% ===== pdf idx 245, printed 231.0 =====
% !PRINTED 231
\weberpage{231}
%<<CONT
\[ \sum \bigl(H\cdot W,\ \Phi(B_i)\bigr)^{\varrho_i} \]
in (3) auf eine Summe von symbolischen Producten reducirt, die nur mehr
durch Faltung von Producten der $B$, also durch Faltung einer Grösse
$\bar V$ entstehen. Dies liefert aber nach (2):
\[ \bar V = \Phi(B_i) + \sum G\cdot W'', \]
oder, weil die Formen $B_i$ zu den Formen $C_i$ gehören:
\[ \bar V = \Phi(C_i) + \sum G\cdot W''. \]

Somit ist auch dieser Term $\equiv \Phi(C_i)$ mod. $G$ und damit die
ganze rechte Seite der Gleichung (3), was zu beweisen war.

§ 21. \gesperrt{Beweis für die Endlichkeit des Formensystems einer
binären Form.}

205. \textit{Gedankengang des Beweises.} Der soeben bewiesene Satz
bildet nun die Grundlage des von Gordan zum ersten Male im 69. Bande des
Borchardt'schen Journals bewiesenen Fundamentalsatzes:

\glqq Jede binäre Form besitzt ein endliches System von Covarianten,
durch welche sich alle unendlich vielen übrigen rational und ganz
darstellen lassen.\grqq{}

Ehe wir hier zum Beweise schreiten, will ich versuchen, dessen
Gedankengang klar zu legen.

Denken wir uns sämmtliche Covarianten von $f=a_x^n$ gefunden und in den
Symbolen $a,\ b,\ c\ldots$ der Form $f$ dargestellt. Jedes dieser
symbolischen Producte können wir nach § 1, Nr. 11 so umformen, dass die
in ihm zumeist auftretende Potenz eines Klammerfactors, etwa des
Factors $(ab)$, eine gerade ist. Ist dann $g$ die grösste in
$\dfrac{n}{2}$ enthaltene ganze Zahl, so können wir sämmtliche
Covarianten von $f$ in folgende $g$ Klassen eintheilen, wobei die
späteren Klassen alle Covarianten der früheren umfassen. Die erste
Klasse enthält alle Covarianten mit dem zumeist auftretenden
Klammerfactor $(ab)^0$; wir bezeichnen sie als System $(A^{(0)})$ und
erkennen, dass es nur aus der Originalform $f=a_x^n$ allein bestehen
kann.

Die zweite Klasse enthält alle Covarianten, deren zumeist auftretender
Klammerfactor höchstens $(ab)^2$ ist; wir bezeichnen sie als System
$(A^{(1)})$.

Das System $(A^{(2)})$ wird sodann alle Covarianten umfassen, deren
zumeist auftretender Klammerfactor höchstens $(ab)^4$ ist etc., und
endlich das System $(A^{(g)})$ alle Covarianten von $f$ überhaupt, da ja
keine
%CONT>>


% ===== pdf idx 246, printed 232.0 =====
% !PRINTED 232
\weberpage{232}
%<<CONT
Covariante einen höheren Klammerfactor als $(ab)^{2g}$ besitzen kann,
insoferne eben $2g=n$, oder $n-1$ ist. \gesperrt{Jedes solche System
$(A^{(k)})$ ist relativ vollständig in Bezug auf den Modul $(ab)^{2k+2}$,
und alle höheren Potenzen von $(ab)$.} Denn durch Faltung eines
Productes von Formen $A_i^{(k)}$ können nur symbolische Producte
entstehen, deren zumeist auftretender Klammerfactor $(ab)^{2\varrho}$,
$\varrho \leqq k$, ist, und symbolische Producte mit dem Factor
$(ab)^{2\varrho}$, $\varrho > k$.

Da nun aber das System $(A^{(k)})$ sämmtliche Covarianten enthalten
soll, deren zumeist auftretender Klammerfactor höchstens $(ab)^{2k}$
ist, so kann die erste Sorte von Producten nur wieder auf Formen
$A_i^{(k)}$ führen, und es ist demnach in der That:
\[ \bar U^{(k)} \equiv F(A_i^{(k)}) \ \text{mod. } (ab)^{2k+2},\ (ab)^{2k+4}\ \text{etc.,} \]
und insbesondere
\[ \bar U^{(g)} \equiv F(A_i^{(g)}) \ \text{mod. } (ab)^{2g+1}, \]
wobei $(A^{(g)})$ das volle System von $f=a_x^n$.\footnote{*)
\textit{Beispiel.} Bei Formen vierter Ordnung besteht das System
$(A^{(0)})$ aus der Form
\[ f = a_x^4; \]
das System $(A^{(1)})$ aus den Formen:
\[ f,\ \Delta,\ t,\ j; \]
das System $(A^{(2)})$ aus den Formen:
\[ f,\ \Delta,\ t,\ j,\ i, \]
und es ist:
\[
\begin{aligned}
\bar U^{(0)} &\equiv F(f) + \sum (ab)^2\cdot W\\
\bar U^{(1)} &\equiv F(f,\ \Delta,\ t,\ j) + \sum (ab)^4\cdot W \qquad\text{(vergl. Nr. 202)}\\
\bar U^{(2)} &\equiv F(f,\ \Delta,\ t,\ j,\ i).
\end{aligned}
\]}

Hiebei ist das erste System $(A^{(0)})$ endlich; denn es besteht aus der
Form $f$ allein. Gelingt es uns nun durch Ueberschiebung des Systemes
$(A^{(0)})$ mit einem andern gleichfalls endlichen Systeme $(B^{(0)})$
das System $(A^{(1)})$ zu erhalten, und weiter durch Ueberschiebung des
Systemes $(A^{(1)})$, das nun nothwendig nach Nr. 199 gleichfalls
endlich ist, mit einem endlichen System $(B^{(1)})$ das System
$(A^{(2)})$ zu erzeugen etc. etc., so ist der Beweis der Endlichkeit
geliefert. Solche brauchbare Systeme $(B^{(k)})$ liefert aber sofort der
bereits in Nr. 43 geschilderte Process, und ich gehe nun dazu über,
diese Bildsysteme aufzustellen.

206. \textit{Aufstellung der Bildsysteme $(B^{(k)})$.} Es sei $A_i^{(k)}$
irgend eine Form des Systemes $(A^{(k)})$; sie sei durch eine gewisse
Anzahl Faltungen
%CONT>>


% ===== pdf idx 247, printed 233.0 =====
% !PRINTED 233
\weberpage{233}
%<<CONT
aus einer Potenz $f^\varrho = a_x^n \cdot b_x^n \cdot c_x^n \cdots r_x^n$ gewonnen worden. Nehmen
wir nun \gesperrt{genau dieselben} Faltungen an der $\varrho^{\mathrm{ten}}$ Potenz der Covariante
$\varphi = (ab)^{2k+2} a_x^{n-2k-2} b_x^{n-2k-2}$ vor, was (vergl. Nr. 43) auf mehrere
Weisen möglich ist, so entsteht wiederum eine Covariante von $f$, die
wir mit $B_i^{(k)}$ bezeichnen. Wir sagen: die Form $B_i^{(k)}$ ist nach dem
\gesperrt{Modell} $A_i^{(k)}$ aus der Covariante $\varphi = (f,f)^{2k+2}$ gebildet worden. Auf
dieselbe Weise können wir nach jeder Form des Systemes $(A^{(k)})$ ein
Bild herstellen, so lange der Grad der Covariante $\varphi$ grösser oder
gleich $n$ ist, und wir nennen alsdann die Gesammtheit dieser Formen
ein \gesperrt{System} $(B^{(k)})$. Nur wenn der Grad $2n - 4k - 4$ dieser Form
$\varphi$ kleiner wird als $n$, kann diese Operation unausführbar werden, da
alsdann $\varphi^\varrho$ möglicher Weise nicht mehr die nöthige Anzahl Factoren
erster Art besitzt, durch deren Faltung das Bild $B_i^{(k)}$ nach dem Muster
von $A_i^{(k)}$ entstehen soll. In diesem Falle aber nehmen wir alsdann
die vollen Systeme der Formen niedrigerer Ordnung, also der Formen
$(n-1)^{\mathrm{ten}}$, $(n-2)^{\mathrm{ten}}$ etc. Grades zum Muster für die Bildung des
Systemes $(B^{(k)})$. Da wir die Endlichkeit des Systemes einer Form
$n^{\mathrm{ter}}$ Ordnung beweisen wollen, so dürfen wir die vollen Systeme der
Formen niedrigerer Ordnung als endliche annehmen, wie wir dies ja
auch in der That bei Formen ersten bis vierten Grades gesehen haben.

Um zu wiederholen: Ist $k \leqq \frac{n}{4} - 1$, so stellen wir das System
$(B^{(k)})$ nach dem Modelle der vorhin Nr. 205 definirten Systeme $(A^{(k)})$
an der Form $\varphi = (f,f)^{2k+2}$ her. Alle spätern Systeme $(B^{(k)})$ haben
die vollen Systeme $(\overline{A}^{(k)})$ der Formen niedrigeren Grades zum Muster,
die wir als endliche Systeme voraussetzen. Die Systeme $(A^{(k)})$ von
$k = 0$ bis $k \leqq \frac{n}{4} - 1$ genügen, wie wir in Nr. 205 gesehen, der
Congruenz:
\weq{1}{\overline{U}^{(k)} \equiv F(A_i^{(k)}) \ \mathrm{mod.}\ (ab)^{2k+2},\ (ab)^{2k+4},\ \text{etc.};}
die späteren Systeme $(\overline{A}^{(k)})$ aber, als endliche Systeme, der Gleichung
\weq{2}{\overline{U}^{(k)} = F(\overline{A}_i^{(k)}).}
207. \textit{Eigenschaft der Bildsysteme $(B^{(k)})$.} Während so die Modellsysteme
, wenn sie nicht a priori endlich sind, der Congruenz (1) genügen
, befriedigen sämmtliche Bildsysteme die Congruenz:
\weq{3}{\overline{U}^{k} \equiv F(A_i^{(k)}) \ \mathrm{mod.}\ (ab)^{2k+4},\ (ab)^{2k+6} \ldots} % checked against the GDZ scan (27 Oct 2026): under the ink the printed symbol is \overline{U}^{k} (italic U with its rounded bowl and thin right stem, exponent k without parentheses; the printed V of p. 234 is pointed); a reader of this copy wrote V over the U and B over the A of F(A_i^{(k)}), i.e. \overline{V}^{(k)} \equiv F(B_i^{(k)}) as on p. 234; reproduced as printed
Dies erkennt man für die Systeme $(B^{(k)})$, $k > \frac{n}{4} - 1$, unmittelbar.
Denn ihre Modellsysteme genügen der Gleichung (2), d. h. zwischen
%CONT>>


% ===== pdf idx 248, printed 234.0 =====
% !PRINTED 234
\weberpage{234}
%<<CONT
den Formen des Systemes und einer Form $\overline{U}^{(k)}$, die aus ihnen durch
Faltung entsteht, besteht die Relation:
\[
\overline{U}^{(k)} - F(A_i^{(k)}) = 0.
\]
Jedes Bild derselben ausgeführt an $\varphi = (f,f)^{2k+2}$ genügt alsdann
nach Nr. 43, (3) der Relation:
\[
\overline{V}^{(k)} - F(B_i^{k}) \equiv 0 \ \mathrm{mod.}\ (ab)^{2k+4},
\]
d. h. es ist:
\[
\overline{V}^{(k)} \equiv F(B_i^{k}) \ \mathrm{mod.}\ (ab)^{2k+4}.
\]
Für die übrigen Systeme $(B^{(k)})$ von $k = 0$ bis $k \leqq \frac{n}{4} - 1$ ergiebt
 sich die Richtigkeit unserer Behauptung auf folgendem Wege.
Wir schreiben wieder die Congruenz (1), welcher ihre Modelle genügen,
in der Gleichungsform:
\weq{4}{\overline{U}^{(k)} - F(A_i^{(k)}) - \sum{}'(ab)^{2k+2} W - \sum{}'(ab)^{2k+4} W' - \cdots = 0,}
und erkennen in ihr eine Relation von symbolischen Producten, deren
Glieder sich weder gegenseitig aufheben, noch einzeln verschwinden.
Von einer solchen Relation haben wir in Nr. 43 gezeigt, dass ihr
Bild hergestellt an der Form $\varphi = (f,f)^{2k+2}$ sich darstellen lässt durch
symbolische Producte, die mindestens $(ab)^{2k+4}$ als Klammerfactor besitzen
.

Nun ist das Bild von $\overline{U}^{(k)}$ die Form $\overline{V}^{(k)}$, das Bild von $F(A_i^{(k)})$
die Function $F(B_i^{(k)})$; dabei seien die Formen $B_i^{(k)}$, die ja auf verschiedene
 Arten nach dem Modell $A_i^{(k)}$ hergestellt werden können, so
gewählt, dass die Bilder des dritten Termes $\sum{}'(ab)^{2k+2} W$, und auch
der späteren, die Form annehmen
\weq{5}{\sum{}'(ab)^{2k+2}(aa_1)^{2k+2}\overline{W},\ \text{etc.};}
dies ist immer möglich, weil ebenso wie der dritte Term in (4) durch
mindestens $(2k+2)$-malige Faltung aus $f^\varrho = a_x^n b_x^n c_x^n \cdots r_x^n$, auch das
Bild durch mindestens $(2k+2)$-malige Faltung aus
\[
\varphi^\varrho = (ab)^{2k+2} a_x^{n-2k-2} b_x^{n-2k-2} \cdot (a_1 b_1)^{2k+2} a_{1x}^{n-2k-2} b_{1x}^{n-2k-2}
\]
entstehen muss. Allein alle symbolischen Producte, wie sie in der
Summe (5) auftreten, lassen sich durch andere Producte darstellen,
die mindestens den Factor $(ab)^{2k+4}$ besitzen, wie wir in Nr. 98 ausführlich
 gezeigt haben. Es ist also für $k \leqq \frac{n}{4} - 1$ zunächst:
%CONT>>


% ===== pdf idx 249, printed 235.0 =====
% !PRINTED 235
\weberpage{235}
%<<CONT
\[
\begin{gathered}
\overline{V}^{(k)} - F(B_i^{(k)}) - \sum{}'(ab)^{2k+2}(aa_1)^{2k+2}\overline{W} - \sum{}'(ab)^{2k+4}(aa_1)^{2k+4}\overline{W}' \cdots \\
\equiv 0 \ \mathrm{mod.}\ (ab)^{2k+4},
\end{gathered}
\]
oder weil vom dritten Terme links incl. $ab$ alle Producte den nämlichen
 Modul $(ab)^{2k+4}$ als Factor besitzen:
\[
\overline{V}^{(k)} \equiv F(B_i^{(k)}) \ \mathrm{mod.}\ (ab)^{2k+4},\ (ab)^{2k+6}\ \text{etc.}
\]
208. \textit{Rückblick.} Wir haben so die Gesammtheit aller Covarianten
von $f$ in doppelter Weise eingetheilt: Einmal in Systeme $(A^{(k)})$ relativ
vollständig modulo $(ab)^{2k+2}$, $(ab)^{2k+4}$, etc. Dann in Systeme $(B^{(k)})$,
welche die gleiche Anzahl von Formen enthalten, wie das
entsprechende Modell $(A^{(k)})$, welche ferner relativ vollständig sind
modulo $(ab)^{2k+4}$, $(ab)^{2k+6} \ldots$ etc., und welche endlich jedesmal mindestens
eine Form enthalten, die den niedrigsten Modul $(ab)^{2k+2}$ des Modellsystemes
 $A^{(k)}$ zum Factor hat. Denn $(B^{(k)})$ hat immer die jeweilige
Form $\varphi = (f,f)^{2k+2}$ als Originalform. Je zwei Systeme $(A^{(k)})$ und
$(B^{(k)})$ sind also gerade zwei Systeme, wie sie der vierte Lehrsatz
Nr. 204 voraussetzt und durch Ueberschiebung derselben entsteht
sonach ein relativ vollständiges Covariantensystem von $f$, das als
niedrigsten Modul den Modul $(ab)^{2k+4}$ besitzt, d. h. das System
$(A^{(k+1)})$ der Form $f$. Und darin liegt der Hauptschritt des Beweises für
den an die Spitze dieses Paragraphen gestellten Lehrsatzes.

209. \textit{Beweis der Endlichkeit des Formensystemes.} Nachdem wir
nun gesehen haben, wie durch die Hilfssysteme $(B^{(k)})$ aus einem System
$(A^{(k)})$ das nächst höhere System $(A^{(k+1)})$ vermittelst Ueberschiebung
erzeugt wird, erübrigt uns nur mehr, dass wir uns von der Endlichkeit
 der einzelnen relativ vollständigen Systeme $(A^{(k)})$ überzeugen.
Daraus folgt dann von selbst die Endlichkeit des Systemes $(A^{(\varrho)})$,
d. i. des vollen Systemes von $f$. Nun besteht aber das System $(A^{(0)})$
aus der Form $f$ allein und demnach das Bildsystem $(B^{(0)})$ aus
$\varphi = (ab)^2 a_x^{n-2} b_x^{n-2}$ allein; durch ihre Ueberschiebungen entsteht das
System $(A^{(1)})$, das somit endlich ist, weil es die Systeme $(A^{(0)})$ und
$(B^{(0)})$ waren. Nach dem Muster $(A^{(1)})$ bilden wir $(B^{(1)})$; es enthält
dieselbe endliche Anzahl von Formen wie $(A^{(1)})$ und durch Ueberschiebung
 beider entsteht das endliche System $(A^{(2)})$ etc. etc. So gelangen
 wir schliesslich zum vollen System $(A^{(\varrho)})$, von dem wir nun
erkennen, dass es in der That endlich sein muss. In diesem Systeme
$(A^{(\varrho)})$ treten freilich im Allgemeinen noch viele überflüssige Formen
auf, d. h. solche, die sich rational und ganz durch die übrigen des
Systemes ausdrücken lassen. Dieselben zu erkennen bleibt immer eine
der Hauptaufgaben bei der wirklichen Systembildung. Sollte eine
%CONT>>


% ===== pdf idx 250, printed 236.0 =====
% !PRINTED 236
\weberpage{236}
%<<CONT
Form $f = a_x^n$ so beschaffen sein, dass $(f,f)^{2k+2+2\varrho}$, $\varrho = 1, 2, \ldots$
identisch verschwindet, dann bildet natürlich bereits das System
$(A^{(k+1)})$ das vollständige System von $f$. (Vergl. die Systeme der
regulären Körper \S\ 19).*)

\begin{center}
\textbf{\S\ 22. Das Formensystem der Form fünfter Ordnung.}
\end{center}

210. \textit{Aufstellung der Systeme $(A^{(k)})$ und $(B^{(k)})$.} Die eben entwickelte
 Theorie wollen wir nunmehr durch die Aufstellung des Formensystemes
 einer Form fünfter Ordnung erläutern. Das Formensystem
$A^{(0)}$ ist dargestellt durch die Form
\[
(A^{(0)}) \qquad f = a_x^5;
\]
das Formensystem $(B^{(0)})$ durch die Form
\[
(B^{(0)}) \qquad \varphi = (ab)^2 a_x^3 b_x^3 = (f,f)^2.
\]
Die Ueberschiebung beider Systeme liefert das System $(A^{(1)})$. Die
Formen dieses Systemes sind also enthalten in dem Ausdruck
\[
U = \left(f^\alpha,\ ((f,f)^2)^\beta\right)^\lambda.
\]
Es handelt sich nun zunächst darum, jene Ueberschiebungen auszuscheiden
, welche zerfallende oder reducible Glieder enthalten. Hieher
gehören in erster Linie alle jene, für welche $\lambda > 1$ ist, welche Werthe
auch $\alpha > 0$, $\beta > 0$ haben mögen. Denn diese Ueberschiebungen enthalten
 Glieder mit den Klammerfactoren $(ab)^2 (ac)^2$. Solche symbolische
Producte sind aber bereits (nach Nr. 98) $\equiv 0\ \mathrm{mod.}\ (ab)^4$.

$\lambda$ kann also nur den Werth 1 besitzen; in diesem Falle darf aber
wiederum $\alpha$ oder $\beta$ nicht grösser als 1 sein, da ja alsdann schon $f$
und $\varphi$ allein den zur einmaligen Faltung nöthigen Factor erster
Art liefern, während die übrigen Factoren $f$ und $\varphi$ nur multiplicativ
zum Faltungsglied hinzutreten und so dieses Glied zu einem zerfallenden
 Gliede machen würden.

Die einzige neue Form, welche sonach durch die Ueberschiebung
der Systeme $(A^{(0)})$ und $(B^{(0)})$ entsteht, ist die Form
%CONT>>

\bigskip
\noindent\rule{4cm}{0.4pt}

\noindent *) Im 100. Bande des Crelle'schen Journals hat neuerdings auch Herr
\gesperrt{Mertens} einen Beweis für die Endlichkeit des Formensystems erbracht.
Indem er voraussetzt, dass das System $S$ der Form $p = a_x^n$ endlich
sei, bildet er zunächst das simultane System von $p$ und einer linearen
Form $\alpha_{1x}$. Mit Hilfe von diophantischen Gleichungen zeigt er, dass
dieses System eine endliche Anzahl von Formen $B$ enthält, welche
in den Coefficienten von $p$ und $\alpha_1$ vom gleichen Grade sind.
Nun identificirt er die Form $f = A_x^{n+1}$ mit $p \cdot \alpha_1 = p_x^n \cdot \alpha_{1x}$, indem
er $\alpha_1$ als den einen und $p_x^n$ als das Product $\alpha_{2x}\alpha_{3x}\cdots\alpha_{n+1,x}$ der
übrigen linearen Factoren von $f$ betrachtet.


% ===== pdf idx 251, printed 237.0 =====
% !PRINTED 237
\weberpage{237}
\[
t = (f,\, (f,f)^2) = (ab)^2(ca)\, a_x^2 b_x^3 c_x^4
\]
und das System $(A^{(1)})$ besteht somit aus den drei Formen
\[
f,\ (f,f)^2,\ (f,\,(f,f)^2).
\]
Dies galt schon bei den Formen dritten und viertes Grades und % [sic]
gilt auch noch bei h\"oheren Formen. F\"ur die regul\"aren K\"orper ist
das System $(A^{(1)})$, abgesehen von einer Invariante, das volle System.

Nach diesem Modelle w\"are nun an der Form
\[
i = (ab)^4 a_x b_x = (f,f)^4
\]
das System $(B^{(1)}) \equiv 0 \bmod. (ab)^5$ zu bilden. Da jedoch $i$ bereits eine
Form zweiter Ordnung ist, so muss man das System einer Form
$f = a_x^2$ zum Muster nehmen, das nur aus den Formen $a_x^2$ und $(ab)^2$
besteht. Demnach enth\"alt $(B^{(1)})$ die beiden Formen
\[
(B^{(1)}) \qquad i,\ (i,i)^2 = (ab)^4(cd)^4(ac)(bd) = A.
\]
Dieses System ist $\equiv 0 \bmod. (ab)^5$. Durch Ueberschiebung desselben
 mit $(A^{(1)})$ erh\"alt man somit das vollst\"andige System $(A^{(2)})$ der
Form $f = a_x^5$.

211. \textit{Ueberschiebung der Systeme $(A^{(1)})$ und $(B^{(1)})$.} Unsere n\"achste
Aufgabe besteht nun wiederum darin, aus den Ueberschiebungen
$(U,V)^\lambda$ jene auszuscheiden, welche zerfallende Glieder enthalten. Zu
dem Zwecke untersuchen wir zuerst die Ueberschiebungen einzelner
Formen aus beiden Systemen, sodann die Ueberschiebungen von Producten
 $U$ mit Producten $V$. Wir beginnen mit den Ueberschiebungen:
\[
(f^\mu,\, i^\varrho)^\lambda.
\]
Ist hier $\mu = 1$, so ergeben sich von $\lambda = 1$ bis $\lambda = 5$ Ueberschiebungen
 ohne zerfallende Glieder, n\"amlich die Formen
\[
(f,i),\ (f,i)^2,\ (f,i^2)^3,\ (f,i^2)^4,\ (f,i^3)^5,
\]
wovon man sich leicht \"uberzeugt, wenn man das Product
\[
f \cdot i^\varrho
\]
bezw. 1, 2, 3, 4, 5-mal auf alle m\"oglichen Arten faltet.

Bei sechsmaliger Faltung des Productes
\[
f^2 \cdot i^3 = c_x^5 \cdot d_x^5 \cdot (ab)^4 a_x b_x \cdot (a_1 b_1)^4 a_{1x} b_{1x} \cdot (a_2 b_2)^4 a_{2x} b_{2x}
\]
entstehen dagegen zerfallende Glieder; so z.\,B. wenn man die vier ersten
Factoren $c_x$ mit den vier Factoren $a_x b_x a_{1x} b_{1x}$, und zwei Factoren $d_x$
mit $a_{2x} b_{2x}$ faltet. Das so entstehende Glied der Ueberschiebung
$(f^2, i^3)^6$ ist nichts anderes als das Product zweier Glieder niedriger
Ueberschiebungen, n\"amlich der Ueberschiebungen
\[
(f, i^2)^4 \ \text{und}\ (f, i)^2.
\]


% ===== pdf idx 252, printed 238.0 =====
% !PRINTED 238
\weberpage{238}
Ganz ebenso zeigt man, dass die Ueberschiebungen
\[
(f^2, i^4)^7,\ (f^2, i^4)^8,\ (f^2, i^5)^9
\]
zerfallende Glieder enthalten. Dagegen liefert die Ueberschiebung
\[
(f^2, i^5)^{10},
\]
wie man leicht wieder durch den Faltungsprocess erkennt, keine zerfallenden
 Glieder; sie ist eine Invariante vom Grade 12 in den Coefficienten
 von $f$.

Wenn man nun weitere Potenzen von $i$ mit $f^3, f^4, \ldots$ etc. zur
Faltung bringt, so sieht man, dass immer Glieder hiebei gebildet
werden k\"onnen, welche diese Invariante zum Factor haben. Wir erhalten
 sonach von dieser Seite aus keine neuen Formen mehr.

\textit{Die Ueberschiebungen $(\varphi^\mu, i^\varrho)^\lambda$.} Nicht zerfallende Glieder enthalten
 die Ueberschiebungen:
\[
(\varphi, i),\ (\varphi, i)^2,\ (\varphi, i^2)^3,\ (\varphi, i^2)^4,\ (\varphi, i^3)^5,\ (\varphi, i^3)^6,
\]
wie man wieder unmittelbar aus den entsprechenden Faltungsprocessen
erkennt. Hiebei liefert die letzte Ueberschiebung eine Invariante vom
Grade 8 in den Coefficienten von $f$.

Bei allen weitern Ueberschiebungen von Potenzen der Form $\varphi$
mit Potenzen von $i$ kann man sich stets durch Faltung Glieder verschaffen
, welche diese Invariante zum Factor haben, so dass die Reihe
der Formen, die hier entstehen konnten, ebenfalls abgeschlossen ist.

\textit{Die Ueberschiebungen $(t^\mu, i^\varrho)^\lambda$.} Hier enthalten alle ungeraden
Ueberschiebungen von $t$ \"uber $i, i^2, i^3, i^4, i^5$, ausgenommen die neunte
Ueberschiebung, stets reducible Glieder.

Denn setzen wir einen Augenblick
\[
i = (ab)^4 a_x b_x = i_x^2 \qquad
t = (cd)^2 (ce)\, c_x^2 d_x^3 e_x^4,
\]
so tritt bei diesen ungeraden Ueberschiebungen stets ein Glied auf
mit den Faltungen
\[
(ci)(ce)\, i_x c_x.
\]
Dieser Theil des Gliedes kann aber nach dem Productsatz in die
Form gebracht werden:
\[
(ci)(ce)\, i_x c_x = \frac{1}{2} \left\{ (ci)^2 e_x + (ce)^2 i_x^2 - (ie)^2 c_x^2 \right\},
\]
und wenn man daher die rechte Seite f\"ur die linke in das betreffende
Glied substituirt, so spaltet es sich in drei zerfallende Glieder.


% ===== pdf idx 253, printed 239.0 =====
% !PRINTED 239
\weberpage{239}
Dagegen tritt bei den geraden Ueberschiebungen kein solches
Glied auf, weil hiebei jedesmal alle Factoren erster Art $i_x$, und
keines bei der neunten Ueberschiebung, weil hier alle Factoren erster
Art $t_x$ aufgebraucht werden.

Wir erhalten sonach aus den Ueberschiebungen mit $t_x^9$ die Formen
\[
(t, i)^2,\ (t, i^2)^4,\ (t, i^3)^6,\ (t, i^4)^8,\ (t, i^5)^9.
\]
Weitere Ueberschiebungen brauchen hier nicht mehr untersucht
zu werden, da das Quadrat der Form $t$, welche ja Functionaldeterminante
 ist, durch niedrigere Formen sich ausdr\"ucken l\"asst.

212. \textit{Die Ueberschiebungen von Producten von Formen.} Es er\"ubrigt
also nur noch die Ueberschiebungen der Form $i^\varrho$ \"uber das Product
der Formen
\[
f^\alpha \cdot \varphi^\beta \cdot t^\gamma
\]
zu studiren.

Hiebei kann aber, da $i$ quadratisch ist, also bei Faltung sich
seine Linearfactoren h\"ochstens auf zwei Formen vertheilen lassen,
stets einer der Exponenten gleich null genommen werden. Von den
sich so ergebenden drei verschiedenen Producten
\[
t^\gamma \cdot \varphi^\beta,\ \ f^\alpha \cdot \varphi^\beta,\ \ f^\alpha \cdot t^\gamma
\]
k\"onnen die beiden ersten unber\"ucksichtigt bleiben, da sie je eine Form
geraden und je eine ungeraden Grades enthalten und somit stets die
Linearfactoren von $i^\varrho$ bei Faltung zun\"achst auf alle Factoren erster
Art von $\varphi$ vertheilt werden k\"onnen, so dass eine Potenz der Invariante
$(\varphi, i^3)^6$ als Factor auftritt.

Sonach bleiben nur noch die Ueberschiebungen
\[
(f^\alpha \cdot t^\gamma,\, i^\varrho)^\lambda
\]
zu studiren. Dabei muss $2\varrho$ sowohl als $\lambda$ den Grad 9 von $t$ \"ubersteigen
, da ja sonst a priori zerfallende Glieder in der Ueberschiebung
auftreten, n\"amlich die bereits in das System aufgenommenen Ueberschiebungen
 von $t$ \"uber $i^\varrho$.

Setzen wir in erster Linie $\alpha = \gamma = 1$. Die Ueberschiebungen
\[
(f \cdot t, i^5)^{10},\ (f \cdot t, i^6)^{11},\ (f \cdot t, i^6)^{12},\ (f \cdot t, i^7)^{13}
\]
enthalten Glieder, welche in Producte von Gliedern zerfallen, die bereits
betrachteten Ueberschiebungen angeh\"oren. Um nur ein Beispiel anzuf
\"uhren: die Ueberschiebung $(f \cdot t, i^6)^{11}$ enth\"alt ein Glied, in welchem
die sechs Factoren erster Art von $i^3$ mit sechs Factoren von $t$ zur
Faltung kommen, w\"ahrend von den \"ubrigen sechs Factoren f\"unf mit
$f$ gefaltet werden. Dieses Glied ist also ein Product zweier Glieder,
%CONT>>


% ===== pdf idx 254, printed 240.0 =====
% !PRINTED 240
\weberpage{240}
%<<CONT
welche den Ueberschiebungen
\[
(t, i^3)^6,\ \text{bezw.}\ (f, i^2)^5
\]
angeh\"oren.

Dagegen liefert die Ueberschiebung
\[
(f \cdot t, i^7)^{14}
\]
eine neue Invariante, die wir mit $R$ bezeichnen wollen. Ueberschiebungen
 h\"oherer Potenzen von $f$ und $t$ brauchen nicht betrachtet
zu werden; denn jede solche Ueberschiebung enth\"alt immer Glieder,
die $R$ als Factor haben.

213. \textit{Ueberblick \"uber das vollst\"andige Formensystem von $f = a_x^5$.}
Wir haben somit das ganze Formensystem der Form f\"unfter Ordnung
aufgestellt. Es besteht aus den 23 Formen:
\[
\begin{aligned}
1) &\quad f,\ \varphi = (f,f)^2,\ t = (f,\,(f,f)^2),\ i = (f,f)^4,\ A = (i,i)^2 \\
2) &\quad (f,i)^1,\ (f,i)^2,\ (f,i^2)^3,\ (f,i^2)^4,\ (f,i^3)^5 \\
3) &\quad (\varphi,i)^1,\ (\varphi,i)^2,\ (\varphi,i^2)^3,\ (\varphi,i^2)^4,\ (\varphi,i^3)^5,\ (\varphi,i^3)^6 \\
4) &\quad (t,i)^2,\ (t,i^2)^4,\ (t,i^3)^6,\ (t,i^4)^8,\ (t,i^5)^9 \\
5) &\quad (f^2,i^5)^{10},\ (f \cdot t, i^7)^{14}.
\end{aligned}
\]
Es ist freilich noch nicht gezeigt, dass nicht einige dieser Formen
durch die \"ubrigen sich rational und ganz ausdr\"ucken lassen; doch
darf die Annahme gemacht werden, dass dieses System nicht mehr
weiter reducibel ist, da alle bisherigen diesbez\"uglichen Untersuchungen
keine gegentheiligen Resultate zu Tage gef\"ordert haben.

Unter diesen 23 Formen befinden sich drei fundamentale Invarianten,
n\"amlich:
\[
(i,i)^2,\ (\varphi, i^3)^6,\ (f^2, i^5)^{10},
\]
vom beziehungsweise vierten, achten und zw\"olften Grade in den Coefficienten
. Hiezu tritt noch die Invariante
\[
(f \cdot t, i^7)^{14}.
\]
Sie ist vom Grade 18 in den Coefficienten und hat die Eigenschaft
, bei der Transformation
\[
x_1 = y_1,\quad x_2 = -y_2,
\]
deren Modul $\Delta = -1$ ist, das Zeichen zu \"andern. Denn das symbolische
 Product
\[
\bigl(f \cdot (f,\,(f,f)^2),\, i^7\bigr)^{14}
\]
enth\"alt 90 Buchstaben, 70 herr\"uhrend von $i^7$, und viermal je f\"unf
herr\"uhrend von $f$; dieselben vertheilen sich sonach auf 45 Klammerfactoren
, und da bei dieser Transformation jeder solcher Klammerfactor
sein Zeichen \"andert, so \"andert dasselbe die ganze Form, d.\,h. sie ist
eine \glqq schiefe Invariante\grqq{} (vergl. Nr. 8).


% ===== pdf idx 255, printed 241.0 =====
